\documentclass[10pt]{article}
\usepackage{graphicx}
\usepackage{amsmath}
\usepackage{amssymb}
\usepackage{amsfonts}
\usepackage{textcomp}
\usepackage{amsthm}
\usepackage{mathtools}

\theoremstyle{plain}
\newtheorem{theorem}{Theorem}[section]
\newtheorem{lemma}[theorem]{Lemma}
\newtheorem{corollary}[theorem]{Corollary}
\newtheorem{proposition}[theorem]{Proposition}

\theoremstyle{definition}
\newtheorem{definition}[theorem]{Definition}
\newtheorem{example}[theorem]{Example}

\theoremstyle{remark}
\newtheorem{remark}[theorem]{Remark}
\title{%
	\vspace{-1.5em}%
	\large\textbf{Part Eleven}%
}

\author{A. E. Mahrouss}
\date{\today}

\begin{document}
	
	\begin{definition}
		Let the following:
		\begin{align*}
			f'(x)&= \frac{\ln (x)^{\mathrm{e}}}{\mathrm{e}} \cdot \mathrm{e}^{x};\\
			f(x) &= \int_{0}^{1} f'(x)dx;\\
			&= \int_{0}^{1} \frac{\ln (x)^{\mathrm{e}}}{\mathrm{e}} \cdot \mathrm{e}^{x} dx.
		\end{align*}
	\end{definition}
	
	\begin{definition}
		Let the following:
		\begin{align*}
			x &= 2(\sqrt{2} + 1) + \frac{5}{2}.
		\end{align*}
		Thus for $x$:
		\begin{align*}
			\int_{0}^{1} \frac{\ln (x)^{\mathrm{e}}}{\mathrm{e}} \cdot \mathrm{e}^{x} dx.
		\end{align*}
		The following definite integral admits a solution $R$.
		\begin{align*}
			R &= \int_{0}^{1} \frac{\ln (x)^{\mathrm{e}}}{\mathrm{e}} \cdot \mathrm{e}^{x} dx;\\
			&= \int_{0}^{1} \frac{\ln (2(\sqrt{2} + 1) + \frac{5}{2})^{\mathrm{e}}}{\mathrm{e}} \cdot \mathrm{e}^{2(\sqrt{2} + 1) + \frac{5}{2}};\\
			&= \mathrm{e}^{7/2+2 \sqrt{2}} \log^{e}(\frac{9}{2}+2+\sqrt{2}) + C.
		\end{align*}
	\end{definition}
	
	\nocite{*}
	
	\bibliographystyle{acm}
	\bibliography{S0022314X10001836, A_Course_of_Pure_Mathematics, General_Theory_Dirichlet_Series, Riemann}
	
\end{document}